\section{Exterior-power bounds and a positive initial exception}
\label{sec:exterior-witnesses}

\begin{proposition}[PNEW-006: an exterior ambient bound]
Put $g=\gcd(n,3)$, $k=3s/g$, and $d=ns/g$, for $s\ge1$.
Then
\[
 f_n(s)=\dim\operatorname{Hom}_{GL_k}
       ((\det)^n,\bigwedge^d(\bigwedge^3\mathbb C^k))
       \le \binom{\binom{k}{3}}{d}.
 \tag{PNEW-ext}
\]
The binomial coefficient is zero when its lower argument exceeds
its upper argument.  In particular the exact dyadic residue gives
\[
 2W_9(7)=\frac{782264335828}{64}
       >2319959400=\binom{35}{21}\ge f_9(7),
 \tag{PNEW-nine}
\]
so $p(f_9)>7$ and $s(f_9)\ge7$.
There is also a witness with a proved nonzero actual coefficient:
\[
 2W_9(6)=\frac{1819142906755}{64}>1=f_9(6).
 \tag{PNEW-nine-positive}
\]
At rank 15 the certificate gives the larger finite witness
\[
 2W_{15}(60)>2^{881}>2^{866}>\binom{979}{300}\ge f_{15}(60),
 \tag{PNEW-fifteen-sixty}
\]
and hence $p(f_{15})>60$, $s(f_{15})\ge60$.
These finite consequences combine the inherited exact residue with
the exterior identity already proved in ALPOGE\_GCT P-F32.
The contribution here is the explicit comparison and its certificate;
no novelty is claimed for the exterior identity or its top-power case.
\end{proposition}
\begin{proof}
Work first in the ring of symmetric functions over $\mathbb Q$.
Let $\omega$ be the involution $p_r\mapsto(-1)^{r-1}p_r$,
so $\omega(s_\lambda)=s_{\lambda'}$.  The power-sum formulas
\[
 h_3=\frac{p_1^3+3p_1p_2+2p_3}{6},\qquad
 e_3=\frac{p_1^3-3p_1p_2+2p_3}{6}
\]
show directly that
$\omega(p_r[h_3])=(-1)^{r-1}p_r[e_3]$.
Expanding $h_d=\sum_{\alpha\vdash d}p_\alpha/z_\alpha$
and $e_d=\sum_{\alpha\vdash d}(-1)^{d-\ell(\alpha)}
p_\alpha/z_\alpha$ therefore gives
\[
                         \omega(h_d[h_3])=e_d[e_3].
\]
The target partition $(k^n)$ transposes to $(n^k)$.
Specialization to $k$ variables preserves the coefficient of that
Schur function: terms with more than $k$ rows vanish, and the
remaining Schur polynomials are linearly independent.  The last
partition represents $\det^n$ of $GL_k$.  The right hand side is
the character of the exterior power in (PNEW-ext).  Complete
reducibility and its ambient dimension give the inequality.

For $(n,s)=(9,7)$, one has $(k,d)=(7,21)$ and
$\dim\bigwedge^3\mathbb C^7=35$.  The exact residue computation
described below produces the two displayed rational values.  Their
comparison with the integer $\binom{35}{21}=2319959400$ is exact.
For $(n,s)=(9,6)$, put $U=\bigwedge^3\mathbb C^6$.
Then $\dim U=20$ and $\det U=(\det\mathbb C^6)^{10}$,
because each variable occurs in $\binom52=10$ basis triples.
Exterior duality gives
\[
 \bigwedge^{18}U\simeq(\bigwedge^2 U)^*\otimes\det U.
\]
The alternating pairing $U\otimes U\to\bigwedge^6\mathbb C^6$
given by exterior multiplication is nondegenerate.  It is alternating
because three is odd.  The irreducibility of $\bigwedge^3\mathbb C^6$
and Schur's lemma show that this equivariant pairing is unique up
to a scalar.  Hence the multiplicity of $\det$ in $\bigwedge^2U$
is one, and the multiplicity of $\det^9$ in $\bigwedge^{18}U$
is one.  This proves $f_9(6)=1$.
For the rank-15 comparison at $s=60$, use the original symmetric
ambient space: $\dim\operatorname{Sym}^3\mathbb C^{15}=680$ and
$d=15s/3=300$, so complete reducibility gives
$f_{15}(60)\le\dim\operatorname{Sym}^{300}(\mathbb C^{680})
=\binom{979}{300}$.  The two strict power-of-two comparisons in
(PNEW-fifteen-sixty) are verified as integer inequalities in the
certificate below.
\end{proof}

\paragraph{The exact finite computation.}
The script \texttt{scripts/certify\_plethysm.py} uses the positive
multiset $\mathcal S_n$ and performs arithmetic in
$\mathbb Q[z]/(z^{M/2}+1)$ only.  Multiplication by $z^a$ is a
signed permutation of the monomial basis.  Each cycle of that
permutation has total sign $-1$, because $M$ is a power of two
and $a\not\equiv0\pmod M$.  Thus $1-z^a$ is invertible, and the
linear recurrence around each cycle solves
$(1-z^a)v_{\rm new}=2v_{\rm old}$ over the integers.
The script checks that identity after every factor, before
cancelling powers of two.  It independently compares the positive
multiset against the character
$\operatorname{Sym}^3V-V\otimes V^*$ with all multiplicities.
Consequently its computed coefficient is exactly the residue
coefficient in (P-F56v), including the phase.  Formula (P-F56z)
then gives $2W_n(x)$ with no numerical approximation.
The complete rational results and integer inequality margins are
in \texttt{certificates/plethysm\_exact.json}.  Eight witnesses and
all 663 individual inverse-factor identities pass; that count
refers to the four dimensions $4,9,10,15$, not to an extrapolation.

\begin{proposition}[PNEW-007: a nonzero input with strict saturation index at least seven]
Let
\[
 \mu=(3),\qquad\lambda=(35),\qquad\pi=(7^{15}),
 \qquad h(s)=a_{s\lambda,\mu}^{s\pi}=f_{15}(7s).
\]
Then
\[
                      h(1)=1,\qquad p(h)>7,\qquad s(h)\ge7.
 \tag{PNEW-positive-start}
\]
Thus this is an exception inside the nonzero-coefficient population
used in narrow SH(b), rather than merely an input whose initial
coefficient vanishes.
\end{proposition}
\begin{proof}
At $n=15,s=7$, (PNEW-ext) identifies the coefficient with the
multiplicity of $\det^{15}$ in
$\bigwedge^{35}(\bigwedge^3\mathbb C^7)$.  The inner space has
dimension 35, and its determinant is $\det^{\binom62}=\det^{15}$.
The top exterior power is exactly this one-dimensional representation,
so $h(1)=1$.

For $s=14$, the same bound gives
\[
             f_{15}(14)\le\binom{364}{70}<2^{253}.
\]
The exact residue certificate gives
$2W_{15}(14)=A/64$, where $2^{888}<A<2^{889}$.
For completeness, $A$ is the integer obtained by concatenating
the following four digit strings, without spaces or line breaks:
\begin{center}\footnotesize\ttfamily
3333954305645752920064485456028591317103939322607541526148896903832899860\\
4542313254224303373186420275533728219695881859566923177720354053196664566\\
8443320869445013925724828682556583506896582799729547390014645903631244200\\
6703694838860313433573900791230758229801554002547
\end{center}
The script verifies the full cyclotomic recurrence and the strict
integer inequalities, not logarithmic floating-point surrogates.
It follows that $2W_{15}(14)>f_{15}(14)$.

We must check the effect of the odd dilation on constituents, since
an even dilation could discard the primitive dyadic wave.  Let $L$
be a period of $f_{15}$, write $D=\gcd(L,7)$ and $L'=L/D$.
Then $L'$ is a period of $h$, and $v_2(L')=14$.
On residue $a$ modulo $L'$ its constituent is
$q_{7a\bmod L}(7X)$.  Increasing $a$ by $L'/2$ increases
the old residue by $(7/D)L/2$, an odd multiple of $L/2$.
Thus (P-F56d) gives exactly
\[
 q_{7a}(7X)-q_{7(a+L'/2)}(7X)=2W_{15}(7a).
\]
At $a=X=2$ the first value is $h(2)=f_{15}(14)$ and
the second is strictly negative.  No positivity shift at most 2
can make every constituent coefficient nonnegative, and no
saturation shift below 2 can make every nonzero constituent
positive at all positive integers.  This already proves
$p(h)>2$ and $s(h)\ge2$.
For the stronger bounds stated above, the same exact certificate gives
\[
       2W_{15}(49)>2^{875}>2^{766}>
               \binom{924}{245}\ge f_{15}(49).
\]
Here the last bound is the symmetric ambient dimension
$\dim\operatorname{Sym}^{245}(\mathbb C^{680})$.
Apply the identical paired-constituent argument at $a=X=7$.
All displayed power-of-two comparisons are exact integer tests
in the accompanying script.
\end{proof}
